% !Mode:: "TeX:UTF-8"

%? ************************************************************* %?
%?                                                               %?
%?                   您将在此处完成中英文摘要的填写                   %?
%?                                                               %?
%? ************************************************************* %?

% ====================== 摘要 ======================

\chapter*{\texorpdfstring{\centering \heiti \sanhao \textbf{摘\quad 要}}{摘要}}
\addcontentsline{toc}{chapter}{\texorpdfstring{\sihao{摘\qquad 要}}{摘要}}

\linespread{1.5}\selectfont
\fangsong\xiaosi

\vspace{1em}

\xfabstract{

    随着电商物流与智能制造的快速发展，传统仓储搬运方式在效率和人力成本方面面临日益严峻的挑战。为提升仓储搬运自动化水平，本文设计了一款具备举升功能的双轮差速驱动自动导引车（AGV），系统开展了总体方案设计、机械结构设计、电气系统设计、运动学建模、路径规划与循迹控制以及运动控制算法研究。

    在总体方案层面，通过三种驱动方案对比分析，选定双轮差速驱动方案，确定了额定载重$25\,\mathrm{kg}$等关键技术指标。机械层面完成了底盘、升降机构和罩壳设计，对支撑梁进行有限元分析，满载工况下最大应力为$4.891\,\mathrm{MPa}$，远低于Q235A屈服强度$235\,\mathrm{MPa}$，最大位移变形仅为$1.538\times10^{-3}\,\mathrm{mm}$。电气层面以STM32F103C8T6为主控芯片，完成了SMC80S伺服电机与ZJPX115行星减速器选型，设计了TCRT5000红外传感器和HC-SR04超声波传感器的接口电路。

    运动学建模与路径规划层面，推导了双轮差速底盘运动学方程，在MATLAB中对直线、圆弧、S形和8字形四类典型轨迹进行了仿真验证。采用Hybrid A*算法在$10\,\mathrm{m}\times8\,\mathrm{m}$仓储栅格地图上规划了无碰撞路径，路径长度$10.71\,\mathrm{m}$。循迹控制层面，模拟5路TCRT5000红外传感器阵列进行PID循迹仿真，RMS横向误差为$0.8\,\mathrm{cm}$，最大横向误差为$1.6\,\mathrm{cm}$，验证了全局规划与局部循迹协同工作的可行性。

    运动控制算法层面，设计了位姿外环与轮速内环相结合的双环PID控制器，在阶跃响应、定点镇定、直线跟踪和圆形轨迹跟踪四类工况下进行了仿真验证。结果表明：阶跃响应无超调、稳态误差接近零；定点镇定终点位置误差为厘米级、航向误差在$1^\circ$以内；直线跟踪角速度RMSE为$0.0023\,\mathrm{rad/s}$；圆形轨迹跟踪线速度RMSE为$0.0088\,\mathrm{m/s}$、轨迹圆度RMSE为$0.021\,\mathrm{m}$。

    仿真结果表明，本文所设计的举升式AGV在结构、路径规划与运动控制层面均具备可行性，可为后续样机研制和工程应用提供参考。

}

\noindent{\heiti\xiaosi\textbf{关键词：}}
\fangsong{举升式AGV；智能仓储；差速驱动；路径规划；Hybrid A*；PID循迹控制}

% ====================== Abstract ======================

\chapter*{}
\vspace*{-1.15cm}
{\centering\bfseries\sanhao ABSTRACT\par}
\addcontentsline{toc}{chapter}{\sihao{ABSTRACT}}
\vspace{0.5cm}

\linespread{1.5}\selectfont
\rmfamily\xiaosi

\xfenabstract{
    With the rapid growth of e-commerce and intelligent manufacturing, traditional warehouse handling faces increasing challenges in efficiency and labor cost. To improve automation in warehouse material handling, this thesis presents a lifting-type dual-wheel differential-drive Automated Guided Vehicle (AGV), covering overall scheme design, mechanical structure design, electrical system design, kinematic modeling, path planning and line-following control, and motion control algorithms.

    For the overall scheme, the dual-wheel differential drive is selected after comparing three drive configurations, with key specifications including rated payload of $25\,\mathrm{kg}$. For the mechanical structure, the chassis, lifting mechanism, and shell are designed. Finite element analysis of the supporting beam yields a maximum stress of $4.891\,\mathrm{MPa}$, far below the Q235A yield strength of $235\,\mathrm{MPa}$, with a maximum displacement of only $1.538\times10^{-3}\,\mathrm{mm}$. For the electrical system, the STM32F103C8T6 microcontroller serves as the main controller, an SMC80S servo motor is matched with a ZJPX115 planetary reducer, and interface circuits are designed for TCRT5000 infrared and HC-SR04 ultrasonic sensors.

    In kinematic modeling and path planning, the differential-drive chassis kinematic equations are derived, and four typical trajectories (straight, circular, S-shaped, and figure-eight) are simulated in MATLAB. The Hybrid A* algorithm generates a collision-free path of $10.71\,\mathrm{m}$ in a $10\,\mathrm{m}\times8\,\mathrm{m}$ warehouse grid map. In line-following control, PID simulation with a 5-sensor TCRT5000 infrared array achieves an RMS lateral error of $0.8\,\mathrm{cm}$ and a maximum error of $1.6\,\mathrm{cm}$, validating the cooperative feasibility of global planning and local line tracking.

    For motion control, a dual-loop PID controller combining an outer pose loop and an inner wheel-speed loop is designed and validated under four scenarios: step response, point stabilization, straight-line tracking, and circular trajectory tracking. Results show zero steady-state speed error in step response, centimeter-level positioning and sub-$1^\circ$ heading error in point stabilization, an angular velocity RMSE of $0.0023\,\mathrm{rad/s}$ in straight-line tracking, and a velocity RMSE of $0.0088\,\mathrm{m/s}$ with a circularity RMSE of $0.021\,\mathrm{m}$ in circle tracking.

    Simulation results demonstrate that the proposed lifting AGV is feasible in structural design, path planning, and motion control, providing a reference for prototype development and engineering applications.
}

\vspace{1em}

\noindent\textbf{Key words:} lifting AGV; intelligent warehousing; differential drive; path planning; Hybrid A*; PID line-following control

\setcounter{secnumdepth}{2}
